\section{The universal resource-dimension frontier}
\label{dim:section}

Here the technology may be any finite-valued function
$F:\R_+^d\to\R_+$ that is coordinatewise nondecreasing, concave, and
positively homogeneous of degree one. In particular, $F(0)=0$.
Continuity between different support faces is not assumed.
The owners, their eligibility, and their finite coordinatewise endowment
bounds are fixed throughout each mechanism. All real reports in those
boxes, including zero reports, are in its domain.

Write $D_d$ for the supremum, over these technologies and fixed finite
owner models, of the infimal funding factor for nonnegative,
coalition-covering TCWM payments. Define $D_d^{\mathrm{OUT}}$ in the same
way with full coalition outside-production protection. Write $Y_S(b)=\sum_{i\in S}y_i(b)$. Thus coverage means
$Y_S(b)\ge F(\sum_{i\in S}b_i)$ for every coalition, and full OUT means
\begin{equation}\label{dim:out-contract}
 Y_S(b)-Y_S(b-w_S)\ge
 F\left(\sum_{i\in S}w_i\right),
 \qquad 0\le w_i\le b_i\quad(i\in S).
\end{equation}
TCWM replaces the right-hand side by zero. A coalition may contain members
whose reports do not change. A funding factor $K$ means
$\sum_i y_i(b)\le K F(\sum_i b_i)$ at every report, including those with
zero production. These definitions do not presume attainment of a
particular model's infimal factor.

\begin{theorem}[Sharp universal dimension factors]\label{dim:main}
For every integer $d\ge1$,
\begin{equation}\label{dim:frontiers}
 D_d=\frac{d+1}{2},\qquad
 D_d^{\mathrm{OUT}}=1+c(d-1),
 \qquad c\bigl(2-e^{-1/c}\bigr)=1,
\end{equation}
where $c$ is the unique solution in $(1/2,1)$
($c\simeq0.5430633898$).
The upper bounds use explicit payment formulas for arbitrary fixed mixed
owners. At a report with exactly $k\ge1$ active resource coordinates they
give the stronger respective factors $(k+1)/2$ and $1+c(k-1)$; at the zero
report they pay zero. The matching lower bounds already hold for
$F(C)=\min_{1\le j\le d}C_j$ on fixed finite endowment boxes, uniformly
over the number of actual owners.
\end{theorem}

The sharpness is a worst-case statement across finite owner models.
It does not assert the same lower bound for every technology or every fixed
owner count. The lower bounds impose no anonymity, continuity,
measurability, homogeneity, or price representation on the original
mechanisms. The proof below constructs regular prices only as a consequence
of assuming a uniform bound for all finite owner models.

\subsection{Payment-only upper bound and a finite-owner lower bound}

For aggregate report $C=\sum_i b_i$, let $I=\{j:C_j>0\}$ and $k=|I|$.
For $j\in I$ define the independent-uniform face average
\begin{equation}\label{dim:uniform-rule}
 A_j(C)=\int_{[0,1]^{I\setminus\{j\}}}
 F\left(C_j e_j+\sum_{\ell\in I\setminus\{j\}}u_\ell C_\ell e_\ell\right)
 \,du_{-j},\qquad
 y_i(b)=\sum_{j\in I}\frac{b_{ij}}{C_j}A_j(C).
\end{equation}
An integral over an empty set of coordinates is its integrand; all payments
are zero when $I$ is empty.

\begin{proposition}\label{dim:uniform-upper}
The rule \eqref{dim:uniform-rule} is nonnegative, covers every coalition,
is TCWM at all real reports, and has total payment at most
$(k+1)F(C)/2$ when $k\ge1$.
\end{proposition}

\begin{proof}
For coalition aggregate $a$ and fixed outsider aggregate $o$, its
coordinate-$j$ payment is
\[
 \frac{a_j}{a_j+o_j}A_j(a+o).
\]
Both nonnegative factors are nondecreasing under any coordinatewise
increase of $a$. If a new coordinate becomes active, its own contribution
is nonnegative. In each old face average, integrating over that newly
active coordinate cannot lower the old integrand, whose coordinate was
identically zero. This proves finite TCWM, including comparisons between
different support faces.

For coverage, work in the $k$-dimensional active face. A finite concave
function is locally Lipschitz and differentiable almost everywhere on its
relative interior. At such a point $x$, its gradient $g(x)$ is nonnegative,
and homogeneity gives $g(x)\cdot x=F(x)$. Moreover
\begin{equation}\label{dim:gradient-majorant}
 g(x)\cdot z\ge F(z)
 \quad\text{for every nonnegative $z$ supported on this face.}
\end{equation}
This includes boundary $z$: for $0<t<1$, apply the supporting inequality
at $x$ to $(1-t)x+tz$, and combine it with concavity along the same segment.
After subtracting $F(x)$ and dividing by $t$, Euler's identity gives
\eqref{dim:gradient-majorant}, without taking a boundary-continuity limit.

Let $U_j$ be independent uniform random variables. For almost every fixed
$U_{-j}$, one-dimensional absolute continuity on $[\varepsilon,1]$ gives
\[
 C_j\int_\varepsilon^1
 \partial_jF(u_jC_j,U_{-j}C_{-j})\,du_j
 =F(C_j,U_{-j}C_{-j})-F(\varepsilon C_j,U_{-j}C_{-j})
 \le F(C_j,U_{-j}C_{-j}).
\]
The derivatives are nonnegative. Letting $\varepsilon\downarrow0$ and
using Tonelli therefore yields
$\E g_j(U\mathbin{\odot}C)\le A_j(C)/C_j$.
Averaging \eqref{dim:gradient-majorant} proves
$\sum_jz_jA_j(C)/C_j\ge F(z)$. Taking $z$ to be the aggregate of any
actual coalition proves coverage.

Finally, two applications of concavity give
\begin{align*}
 \sum_{j\in I}A_j(C)
 &\le\sum_{j\in I}
 F\left(C_je_j+\tfrac12\sum_{\ell\in I\setminus\{j\}}C_\ell e_\ell\right)\\
 &\le kF\left(\frac{k+1}{2k}C\right)
 =\frac{k+1}{2}F(C).
\end{align*}
For $k=1$ the same formula pays exactly $F(C)$; the empty face is immediate.
\end{proof}

\begin{proposition}[A finite obstruction]\label{dim:finite-lower}
Let $d\ge2$ and $s\ge2$ be integers. On $n=(d-1)s+1$ fixed owners,
each eligible for every resource and each with endowment box $[0,1]^d$,
the technology $F(C)=\min_jC_j$ forces every nonnegative,
coalition-covering TCWM rule to have funding factor at least
\begin{equation}\label{dim:finite-bound}
 1+\frac{(s-1)((d-1)s-1)}{2s^2+(d-3)s-1}.
\end{equation}
\end{proposition}

\begin{proof}
Suppose the factor is $\kappa$, and put $\delta=\kappa-1\ge0$.
Average the rule over all permutations of the owner labels and resource
labels. The fixed domain and technology are invariant under these finite
groups. Nonnegativity, every coverage and TCWM inequality, and the funding
bound are linear and hence are preserved. It suffices to consider this
invariant rule; symmetry is not an assumption on the original mechanism.

At profile $P_0$, put $s$ pure unit reporters in each resource group
$1,\ldots,d-2$, one pure unit reporter in each of the last two groups,
and leave the other $s-1$ actual owners at zero reports. Let the respective
per-owner payments be $a,b,w$, where $w\ge0$. Selecting one owner from
each resource group produces one, and grand production is one. Thus
\[
 (d-2)a+2b\ge1,\qquad
 (d-2)sa+2b+(s-1)w\le\kappa.
\]
Subtracting gives $(d-2)(s-1)a+(s-1)w\le\delta$, whence
\begin{equation}\label{dim:b-bound}
 2b\ge1-\frac{\delta}{s-1}.
\end{equation}
When $d=2$ the terms involving the absent first groups are zero.

At $P_1$, the $s-1$ zero reporters restore a pure unit of resource $d-1$.
TCWM applied to these owners together with the original resource-$(d-1)$
owner shows that this group's new payment is at least $b+(s-1)w\ge b$.
The first $d-1$ groups are now symmetric. If $S$ is their union, then
$Y_S(P_1)\ge(d-1)b$.

At $P_2$, every member of this same coalition $S$ restores its report to
one in each of resources $1,\ldots,d-1$ and zero in resource $d$.
These are feasible increases within the original unit boxes. Leave the
last-resource owner unchanged. Write $m=(d-1)s$. All $m$ members of $S$
now have the same payment $p$, and write $q$ for the last owner's payment.
One member of $S$ together with that owner produces one, so
$p+q\ge1$ while $mp+q\le\kappa$. Consequently
\[
 (d-1)b\le Y_S(P_1)\le Y_S(P_2)=mp
 \le\frac{m}{m-1}\delta.
\]
Combining this with \eqref{dim:b-bound} and rearranging gives
\eqref{dim:finite-bound}. All denominators are positive. All $n$ owners
remain present throughout; in particular, the proof never assumes that a
zero reporter receives zero at a positive-output profile.
\end{proof}

Letting $s\to\infty$ in \eqref{dim:finite-bound} and using
Proposition~\ref{dim:uniform-upper} proves $D_d=(d+1)/2$ for $d\ge2$.
For $d=1$, homogeneity makes $F$ linear on its ray; linear payments and
grand-coalition coverage give $D_1=1$.

\subsection{The logarithmic rule for full outside production}

Put $t_0=1-e^{-1/c}$. The defining equation for $c$ gives
$c(1+t_0)=1$. The function $h(x)=x(2-e^{-1/x})$ has
\[
 h'(x)=2-e^{-1/x}(1+1/x)>1\quad(x>0),
\]
and $h(1/2)<1<h(1)$, proving the asserted existence and uniqueness.
Let $U$ have density
\begin{equation}\label{dim:log-law}
 f_U(u)=\frac{c}{1-u}\ind_{(0,t_0)}(u).
\end{equation}
It has mass $-c\log(1-t_0)=1$ and mean
$c[-t_0-\log(1-t_0)]=c$.
Use one common scalar within each face average:
\begin{equation}\label{dim:log-rule}
 H_j(C)=\E F\left(C_je_j+U\sum_{\ell\ne j}C_\ell e_\ell\right),
 \qquad y_i(b)=\sum_{j:C_j>0}\frac{b_{ij}}{C_j}H_j(C).
\end{equation}
Absent-coordinate terms are zero. The independent-uniform and common-scalar
averages are different rules.

\begin{lemma}[Two scalar inequalities]\label{dim:scalar}
For positive $x\in\R_{++}^k$, $k\ge2$,
\begin{equation}\label{dim:ray-lemma}
 \sum_{j=1}^k\frac{\E\min\{x_j,U\min_{\ell\ne j}x_\ell\}}{x_j}\ge1.
\end{equation}
If $H:[0,1]\to\R_+$ is nondecreasing and concave and $0\le t\le t_0$,
then, with derivatives taken almost everywhere,
\begin{equation}\label{dim:hinge-lemma}
 \int_0^{t_0}[H(u)-uH'(u)]\,du
 =2\int_0^{t_0}H(u)\,du-t_0H(t_0)
 \ge t_0H(t)-tH(1).
\end{equation}
\end{lemma}

\begin{proof}
For \eqref{dim:ray-lemma}, choose one smallest coordinate $m$, let $S$
be the second-smallest coordinate, counting multiplicity, and put $t=m/S$.
The smallest coordinate contributes
\[
 q(t)=\E\min\{1,U/t\}
 =\begin{cases}
 1-c-c(1-t)\log(1-t)/t,&0<t\le t_0,\\
 c/t,&t_0\le t\le1.
 \end{cases}
\]
A second-smallest coordinate contributes $ct$, and all other terms are
nonnegative. If $t\le t_0$, then
$q(t)+ct-1=c(1-t)[-1-\log(1-t)/t]\ge0$.
If $t\ge t_0$, then $q(t)+ct=c(t+t^{-1})\ge2c>1$.

For \eqref{dim:hinge-lemma}, a possible downward jump at zero may be
filled with the right limit. This does not change the integral and only
strengthens the assertion when $t=0$. A finite nondecreasing concave
function is continuous from the left at one: concavity applied to a segment
ending at one proves this by letting the segment parameter tend to one.
After filling zero, therefore, $H$ is continuous on $[0,1]$.
The integral identity follows by integration by parts from a positive
endpoint and then passage to zero; boundedness gives $uH(u)\to0$.

A continuous piecewise-linear nondecreasing concave function is a
nonnegative constant plus a nonnegative linear function and a positive
linear combination of hinges $\min\{u,s\}$, $0<s\le1$.
For a constant $A$, the residual left side minus right side of the
inequality is $tA$; for $Bu$ it is $tB(1-t_0)$.
For the hinge with parameter $s$, the residual is respectively
\[
 \begin{cases}
 s(t-s),&s\le t,\\
 (s-t)(t_0-s),&t<s<t_0,\\
 t(s-t_0),&s\ge t_0.
 \end{cases}
\]
Each is nonnegative. Linear interpolation on increasingly fine grids
approximates a continuous concave function uniformly by such concave
piecewise-linear functions. The integral identity, involving only values
of $H$, passes the inequality to the limit.
\end{proof}

\begin{proposition}\label{dim:out-upper}
The rule \eqref{dim:log-rule} satisfies full OUT and coalition coverage at
every real report, and its total payment is at most
$[1+c(k-1)]F(C)$ when $k\ge1$.
\end{proposition}

\begin{proof}
Nonnegativity is immediate. As in the proof of
Proposition~\ref{dim:uniform-upper}, the coalition's coordinate-$j$
payment is the product of its nondecreasing share $a_j/(a_j+o_j)$ and
the nondecreasing face value $H_j(a+o)$. The same formula also shows that
an activated coordinate cannot lower any old contribution. Thus finite
TCWM holds, including across support faces. Two Jensen inequalities give
\[
 \sum_jH_j(C)
 \le\sum_jF\left(C_je_j+c\sum_{\ell\ne j}C_\ell e_\ell\right)
 \le[1+c(k-1)]F(C).
\]

First suppose $F$ is smooth on the positive relative interior of every
face. Fix a positive aggregate $C$ in its active face, coalition holdings
$0\le a\le C$, and a nonnegative restoration direction $z$ on this face.
Write $\alpha_j=a_j/C_j$ and
$v_j(u)=C_je_j+u\sum_{\ell\ne j}C_\ell e_\ell$.
Differentiating while both $a$ and $C$ increase in direction $z$ gives
\begin{align}
 D Y_S[z]&=\sum_j\{(1-\alpha_j)L_j(z)+\alpha_j M_j(z)\},
 \label{dim:holding-derivative}\\
 L_j(z)&=\frac{H_j(C)}{C_j}z_j,\nonumber\\
 M_j(z)&=\E\left[U\nabla F(v_j(U))\cdot z
 +(1-U)\partial_jF(v_j(U))z_j\right].\nonumber
\end{align}
The gradients are nonnegative and satisfy
\eqref{dim:gradient-majorant}, so $L_j\ge0$ and $M_j\ge cF(z)$.
The expression is affine in each holding fraction $\alpha_j$.
It is therefore enough to prove $DY_S[z]\ge F(z)$ at every corner
$\alpha\in\{0,1\}^k$.

The result is immediate if $F(z)=0$. For $k=1$ the technology is linear
on the active ray, and the assertion follows directly; assume $k\ge2$.
Initially suppose $z$ is strictly
positive on the active face. If at least two holding fractions are one,
the derivative is at least $2cF(z)\ge F(z)$.
If all are zero, put $x_j=C_j/z_j$. Monotonicity and homogeneity imply
\[
 \frac{F(v_j(u))}{F(z)}
 \ge\min\{x_j,u\min_{\ell\ne j}x_\ell\}.
\]
Equation~\eqref{dim:ray-lemma} yields $\sum_jL_j(z)\ge F(z)$.

It remains to consider a corner with only $\alpha_j=1$.
Normalize $G(x)=F(z\mathbin{\odot}x)/F(z)$, so $G(\mathbf1)=1$.
Common rescaling of $x=C/z$ leaves the normalized directional expressions
in \eqref{dim:holding-derivative} unchanged. We may therefore arrange
$\min_{\ell\ne j}x_\ell=1$. Put $t=x_j$ and $V=G(x)$.
Every face vector is at least $u x$, so the other coordinate terms satisfy
\[
 \sum_{\ell\ne j}\frac{L_\ell(z)}{F(z)}
 \ge cV\sum_{\ell\ne j}\frac1{x_\ell}\ge cV.
\]
The first part of $M_j/F(z)$ is at least $c$.
If $V\ge t_0$, their sum is at least $c(1+t_0)=1$.
If $t\ge1$, monotonicity gives $V\ge1$, so this case is already covered.
The only remaining case is $0<t\le V<t_0$.

Set $H(u)=G(te_j+u\sum_{\ell\ne j}x_\ell e_\ell)$.
Then $H(t)\ge t$, $H(1)=V$, and Euler's identity gives
\[
 \partial_jG(te_j+u x_{-j})
 =\frac{H(u)-uH'(u)}{t}.
\]
Because $(1-u)f_U(u)=c$ on $(0,t_0)$, the second part of $M_j/F(z)$
is at least
\[
 \frac{c}{t}\int_0^{t_0}[H(u)-uH'(u)]\,du
 \ge c(t_0-V)
\]
by \eqref{dim:hinge-lemma}. The singleton corner therefore contributes at
least $c+c(t_0-V)+cV=1$. This completes every holding corner.

Differentiation under the expectation is justified locally by Euler's
identity and nonnegative gradients:
\[
 C_j\partial_jF(v_j(u))\le F(C),\qquad
 uC_\ell\partial_\ell F(v_j(u))\le F(C)\quad(\ell\ne j).
\]
These bounds dominate the differentiated integrands on every compact
positive capacity set, including as $u\downarrow0$.
For directions with zero coordinates, apply the proved inequality to
$z+\eta\mathbf1_I$ and let $\eta\downarrow0$. The derivative expression
is a finite linear function of its direction, while monotonicity gives
$F(z+\eta\mathbf1_I)\ge F(z)$. No continuity of $F$ across faces is used.

For a finite restoration $z$, take $a(s)=a(0)+sz$ and
$C(s)=o+a(s)$, $0\le s\le1$. For every $s>0$ its active face is the
fixed face of $C(1)$. The derivative inequality integrates on
$[\varepsilon,1]$ to
$Y_S(1)-Y_S(\varepsilon)\ge(1-\varepsilon)F(z)$.
Finite TCWM gives $Y_S(\varepsilon)\ge Y_S(0)$ even if the initial
activation produces a payment jump. Letting $\varepsilon\downarrow0$
proves full OUT for smooth $F$. A one-coordinate face is handled by its
linear formula and the zero face is immediate.

Finally, choose a smooth probability density for $Z$ supported in
$[-1,1]^d$ and let
\[
 F_\varepsilon(x)=\E_Z
 F\bigl(\operatorname{diag}(e^{\varepsilon Z_j})x\bigr).
\]
This preserves nonnegativity, monotonicity, concavity, and degree-one
homogeneity. Convolution in logarithmic coordinates makes it smooth on
every positive relative face. Monotonicity and homogeneity give the
pointwise bounds
\begin{equation}\label{dim:smoothing-bound}
 e^{-\varepsilon}F(x)\le F_\varepsilon(x)
 \le e^\varepsilon F(x)
 \qquad\text{on every face.}
\end{equation}
The smooth finite OUT inequalities thus pass to the limit, since
\eqref{dim:smoothing-bound} also bounds each face average and each
nonnegative owner payment multiplicatively. This proves the claim for the
entire stated class, including downward boundary jumps.
Restoring a coalition's entire bundle from its zero report then proves
coverage: its initial payment under \eqref{dim:log-rule} is zero.
\end{proof}

\subsection{Why arbitrary finite-owner rules satisfy the matching OUT lower bound}

The remaining argument concerns only the continuous technology
$F_r(C)=\min_{1\le j\le r}C_j$. Its continuity imposes no restriction on
the broader class used for the upper bound. We give the reduction in full
to distinguish a fixed finite-owner impossibility from a restriction to
price mechanisms.

\begin{lemma}[Prices forced by bounded finite-owner models]
\label{dim:price-reduction}
Fix $r\ge2$. Suppose that, for every $N$, the model with $N$ fixed
pure-resource owners of each of the $r$ types, each with report interval
$[0,1]$, admits a nonnegative, coalition-covering, full-OUT rule funded by
$\kappa F_r$. Then there are nonnegative degree-zero homogeneous prices $p(C)$ on
$\R_{++}^r$, locally Lipschitz there, such that
\begin{align}
 C\cdot p(C)&\le\kappa F_r(C),\label{dim:price-budget}\\
 a\cdot p(C)&\ge F_r(a)\qquad(0\le a\le C),
 \label{dim:price-coverage}\\
 a\cdot p(C)&\ge(a-u)\cdot p(C-u)+F_r(u)
 \qquad(0\le u\le a\le C).
 \label{dim:price-out}
\end{align}
When an aggregate has a zero coordinate, all its prices may be set to
zero and these statements still hold at every nonnegative real profile.
The prices can also be chosen symmetric under resource permutations.
\end{lemma}

\begin{proof}
For each $N$, average its rule over permutations within each resource
type. At a binary profile with $m_j$ active unit reporters of type $j$,
write $\pi_j^{(N)}(m)$ for an active owner's payment. For integer
$0\le u\le a\le m$, coverage, budget, and the OUT comparison that
turns $u_j$ coalition members of type $j$ completely off give
\begin{align}
 m\cdot\pi^{(N)}(m)&\le\kappa F_r(m),\nonumber\\
 a\cdot\pi^{(N)}(m)&\ge F_r(a),\label{dim:count-array}\\
 a\cdot\pi^{(N)}(m)&\ge
 (a-u)\cdot\pi^{(N)}(m-u)+F_r(u).\nonumber
\end{align}
Inactive reporters may have positive payments. We have only dropped their
nonnegative payments from the budget and from the right side of the last
inequality. Set an unused coordinate $\pi_j(m)$ to zero when $m_j=0$.
For $m_j>0$ the coordinate bound is
$0\le\pi_j^{(N)}(m)\le\kappa F_r(m)/m_j$, independently of $N$.
A diagonal subsequence as $N\to\infty$ therefore gives one array
$\pi(m)$ satisfying \eqref{dim:count-array} for all integer count vectors.
Every finite group of these constraints is available once $N$ exceeds
its largest count. The set $\mathcal C$ of arrays obeying these constraints
and coordinate bounds is nonempty, compact in the countable product
topology, and convex.

For a positive integer $h$ let $(T_h\pi)(m)=\pi(hm)$.
Multiplying all counts in \eqref{dim:count-array} by $h$ and dividing by
$h$ shows that $T_h$ preserves $\mathcal C$. These linear maps commute.
For the first $L$ primes $\ell_1,\ldots,\ell_L$, average $T_h\pi$ over
$h=\prod_{j=1}^L\ell_j^{e_j}$, $0\le e_j<L$.
For a fixed prime $\ell_j$, shifting its exponent changes this average
by at most $\kappa F_r(m)/(Lm_v)$ in coordinate $(m,v)$, by cancellation
of the interior terms. A diagonal subsequential limit as $L\to\infty$
is fixed by every prime dilation, hence by every positive integer dilation.

For rational $C>0$ define $p(C)=\pi(LC)$, where $L$ clears the
denominators. Integer-dilation invariance makes this independent of $L$
and homogeneous under positive rational scaling. Clearing the
denominators of $C,a,u$ in \eqref{dim:count-array} proves
\eqref{dim:price-budget}--\eqref{dim:price-out} for rational data.

Regularity follows from finite comparisons, rather than from any property
of the original rules. For rational $C,C'>0$, put
$\lambda=\min_j C_j/C'_j$ and $v=C-\lambda C'\ge0$.
At aggregate $\lambda C'$, give a coalition all of resource $j$ and no
other resource; then restore $v$ to that same coalition. TCWM, which is
implied by OUT, and rational homogeneity yield
\begin{equation}\label{dim:relative-comparison}
 p_j(C')-p_j(C)
 \le\frac{v\cdot p(C)}{\lambda C'_j}.
\end{equation}
On any compact box $[a,b]^r$ with $a>0$, the coordinate prices are bounded
by $\kappa b/a$, the denominator is at least $a^2/b$, and
\[
 0\le v_\ell\le\|C-C'\|_\infty(1+b/a).
\]
Interchanging $C,C'$ supplies the opposite inequality. Thus $p$ is
Lipschitz on the rational points of each such box and has a unique locally
Lipschitz extension to the positive orthant. Its inequalities and
homogeneity pass to real data by continuity. If coalition holdings or
withheld amounts have zero coordinates, approximate retained holdings $a-u$, withdrawals $u$, and outsider
holdings $C-a$ separately by nonnegative rational vectors, retaining their zero
patterns, and reconstruct $a,C$ by summation. This preserves
$0\le u\le a\le C$; both positive aggregate totals remain positive.

If $F_r(C)=0$, every subbundle of $C$ has zero worth, so zero prices give
coverage and budget on that face. If $C>0$ but $F_r(C-u)=0$, the right
payment in \eqref{dim:price-out} is zero, while coverage and monotonicity
give $a\cdot p(C)\ge F_r(a)\ge F_r(u)$.
These observations complete all boundary comparisons. Finally, average
the constructed prices over the finite resource-permutation group.
All the inequalities, homogeneity, and local Lipschitz continuity survive,
and the resulting prices are resource-symmetric.
\end{proof}

\begin{lemma}[Removing abundant coordinates]\label{dim:pair-limit}
Under the hypothesis of Lemma~\ref{dim:price-reduction}, there is a
symmetric, homogeneous, locally Lipschitz price rule for $\min\{A,B\}$
with nonnegative payments, coalition coverage, full OUT, finite funding
factor $\kappa$, and the expenditure cap
\begin{equation}\label{dim:pair-cap}
 B P_2(A,B)\le D A\quad(B\ge A>0),
 \qquad D=\frac{\kappa-1}{r-1},
\end{equation}
together with its symmetric counterpart.
\end{lemma}

\begin{proof}
At $(A,M,\ldots,M)$, $0<A<M$, symmetry gives a short-resource price
$q$ and common long-resource price $b$. Coverage of a bundle with $A$
units of every resource and the budget yield
\[
 q+(r-1)b\ge1,\qquad Aq+(r-1)Mb\le\kappa A.
\]
Subtracting $A$ times the first inequality gives
$Mb\le D A M/(M-A)$.
At $(A,B,M,\ldots,M)$ with $A\le B\le M$, restore a coalition owning
all $B$ units of the second resource to $M$ units. TCWM gives
\begin{equation}\label{dim:prelimit-cap}
 Bp_2(A,B,M,\ldots,M)\le D\frac{AM}{M-A}.
\end{equation}
For $r=2$, the left side is independent of $M$; letting $M\to\infty$
already proves the lemma using the prices from
Lemma~\ref{dim:price-reduction}. Suppose henceforth $r\ge3$.

For each positive rational pair $(A,B)$, the first two prices at
$(A,B,M,\ldots,M)$ lie in $[0,\kappa]$. A common diagonal subsequence
$M\to\infty$ gives rational pair prices $P_1,P_2$.
They inherit symmetry and pair budget at most $\kappa\min\{A,B\}$.
Equation~\eqref{dim:prelimit-cap} gives \eqref{dim:pair-cap}.
For coalition pair holdings $(a,b)\le(A,B)$, add $\min\{a,b\}$ units
of every extra resource to this coalition. The added payment is at most
\[
 \min\{a,b\}\sum_{j\ge3}p_j(A,B,M,\ldots,M)
 \le\frac{\kappa\min\{a,b\}\min\{A,B\}}{M},
\]
which vanishes. Its worth is $\min\{a,b\}$, so pair coverage survives.

For OUT, it is essential to control the shifted extra totals in a
deviation. For fixed $w\ge0$, compare
$z=(A,B,M,\ldots,M)$ and $z'=(A,B,M-w,\ldots,M-w)$, where $M>w$.
Equation~\eqref{dim:relative-comparison}, first in one orientation and
then the other, gives for $j=1,2$
\begin{align}
 p_j(z')-p_j(z)&\le\frac{w\kappa\min\{A,B\}}{Mz_j},
 \label{dim:shift-bound}\\
 p_j(z)-p_j(z')&\le
 \frac{w\kappa\min\{A,B\}}{(M-w)z_j}.\nonumber
\end{align}
Indeed, the forward residual consists only of $w$ units in each extra
coordinate; the reverse scaling is $(M-w)/M$ and its residual lies in
the first two coordinates. The budget bounds both residual expenditures.

Fix rational pair holdings $(a,b)$ and withdrawals $(u,v)\le(a,b)$,
and put $w=\min\{u,v\}$. Give that coalition $w$ units of every extra
resource and withdraw all of them. Its full outside production is $w$.
The high extra-coordinate payment vanishes as above, and its low retained
extra holdings are zero. At the low pair $(A-u,B-v)$ with positive
production, \eqref{dim:shift-bound} shows that the prices with extra totals
$M-w$ have the same subsequential limit as those with extra totals $M$.
If that low pair has zero production, its entire payment is zero by budget.
The limiting inequality is therefore
\begin{align*}
 aP_1(A,B)+bP_2(A,B)
 \ge{}&(a-u)P_1(A-u,B-v)\\
 &+(b-v)P_2(A-u,B-v)+\min\{u,v\}.
\end{align*}
Set prices to zero at zero-production pairs. We have obtained a rational
pair rule satisfying all required finite inequalities and the cap.

This limit need not yet be homogeneous. Consider the compact convex set
of rational pair-price arrays obeying these inequalities, symmetry, the
budget, and \eqref{dim:pair-cap}. Each coordinate is bounded by
$\kappa\min\{A,B\}/A$ or $\kappa\min\{A,B\}/B$.
The operations $(T_hP)(A,B)=P(hA,hB)$ preserve the set. Repeat the
prime-dilation averaging in Lemma~\ref{dim:price-reduction}; all its
constraints are homogeneous and closed, so the resulting rational array
is homogeneous and retains the cap. The finite comparison
\eqref{dim:relative-comparison} now gives local Lipschitz continuity and
extension to real positive pairs exactly as before. Coverage handles any
zero-production low profile as in the last paragraph of that lemma.
\end{proof}

\begin{lemma}[The two-resource obstruction]\label{dim:pair-obstruction}
Every pair-price rule with the properties of
Lemma~\ref{dim:pair-limit} has $D\ge c$.
\end{lemma}

\begin{proof}
For $A\ge B>0$ write $t=B/A$, $p(t)=P_1(A,B)$ and $q(t)=P_2(A,B)$.
These are well defined by homogeneity, locally Lipschitz for $0<t\le1$,
and $p(t)\le Dt$. Budget gives $q(t)\le\kappa$, while symmetry at the
positive balanced profile gives $q(1)=p(1)\le D$.

If a coalition owns all of the short resource and restores only that
resource, TCWM implies that $z(t)=tq(t)$ is nondecreasing. To obtain the
OUT constraint, start at aggregate $(A,B)$ with coalition holdings
$(0,B)$ and restore $(h,h)$ to that same coalition. Its new holdings are
$(h,B+h)$ and its new aggregate is $(A+h,B+h)$, with outsiders fixed.
This is a feasible restoration from a zero holding, not a withdrawal
from that zero holding. At almost every $t<1$, differentiating its finite
OUT inequality at $h=0+$ gives
\[
 p(t)+q(t)+t(1-t)q'(t)\ge1.
\]
Using the cap, this yields
\[
 \left[\frac{tq(t)}{1-t}\right]'
 =\frac{q(t)+t(1-t)q'(t)}{(1-t)^2}
 \ge\frac{1-Dt}{(1-t)^2}.
\]
Integrate from $\varepsilon$ to $t$ and send $\varepsilon\downarrow0$.
Boundedness of $q$ makes the lower endpoint vanish, and hence
\[
 z(t)\ge (1-D)t-D(1-t)\log(1-t).
\]
Since $z$ is nondecreasing and $z(1)=q(1)\le D$, the right side is at
most $D$ for every $0<t<1$. If $D=0$, this is already a contradiction.
For $D>0$, its maximum occurs at $t=1-e^{-1/D}$ and equals
$1-D+De^{-1/D}$. Consequently
$D(2-e^{-1/D})\ge1$. The strict monotonicity of $h$ proved above implies
$D\ge c$.
\end{proof}

\begin{proof}[Completion of Theorem~\ref{dim:main}]
The TCWM equality was proved after Proposition~\ref{dim:finite-lower}.
Proposition~\ref{dim:out-upper} proves the OUT upper bound.
For $d\ge2$, suppose a factor $\kappa<1+c(d-1)$ were feasible for
every fixed finite pure-owner unit-box model of $F_d=\min_jC_j$.
Lemmas~\ref{dim:price-reduction}--\ref{dim:pair-obstruction} would give
$(\kappa-1)/(d-1)\ge c$, a contradiction. Therefore for each such
$\kappa$ at least one finite, fixed, bounded owner model has no
$\kappa$-funded rule. This is a compactness contrapositive over finite
models, not a mechanism on an infinite owner population.
If the supremum of their infimal feasible factors were smaller than the
claimed constant, choosing $\kappa$ strictly between them would give just
such a rule for every model. This also deals with possible nonattainment
of individual infima. For $d=1$, linear payments satisfy full OUT with
factor one and coverage forces that factor.
\end{proof}

\begin{remark}\label{dim:scope}
These dimension frontiers concern protection of total coalition payments,
with the full value of withheld production added only in the OUT contract.
They do not settle a sharp frontier for individual payment monotonicity.
The displayed integrals are mathematical payment formulas; no efficient
evaluation guarantee for an arbitrary value oracle is asserted. Intermediate
outside-protection parameters and extra deletion-marginal floors are not
included in the present theorem.
\end{remark}
